% ============================================================================
%  PRECÁLCULO — Sesión 10
%  Semana 10: 12 al 17 de octubre de 2026   (lunes 12: festivo)
%  2.4 Logaritmos
%  Duración prevista: 3 horas
% ============================================================================
\documentclass[aspectratio=169,10pt]{beamer}

\usepackage[utf8]{inputenc}
\IfFileExists{spanish.ldf}{\usepackage[spanish,es-noshorthands]{babel}}{}
\usepackage{amsmath,amssymb}
\usepackage{booktabs}
\usepackage{tikz}

\usetheme{Uniminuto}

\uniasignatura{Precálculo}
\uniprofesor{Andrés Felipe Puerta Velez}
\unicorreo{andres.puerta-v@uniminuto.edu.co}
\uninrc{MATE1161}
\unisemana{Sesión 10 — Semana 10 · 12 al 17 de octubre de 2026}

\newcommand{\listatiny}{\setbeamerfont{itemize/enumerate body}{size=\tiny}\tiny}

\begin{document}

\portada

\begin{frame}{Orden del día}
  \begin{enumerate}
    \item Llamado a lista y revisión de los ejercicios propuestos
    \item Repaso: dónde quedamos
    \item \azul{Tema 2.4: qué es un logaritmo y de dónde sale}
    \item Forma exponencial y forma logarítmica
    \item Logaritmo decimal, natural y en base 2
    \item Ecuaciones exponenciales simples; taller y cierre
  \end{enumerate}

  \vspace{0.25cm}
  {\footnotesize\color{uniGris} Sesión de 3 horas. El logaritmo es la
  \destacar{tercera operación inversa} del curso: la resta deshace la suma, la
  división deshace el producto, y el logaritmo deshace la potencia.}
\end{frame}

\begin{frame}{Repaso: dónde quedamos}
  \scriptsize
  \begin{block}{Las tres preguntas que se pueden hacer sobre $2^{3}=8$}
    \scriptsize
    \begin{center}
    \begin{tabular}{@{}p{3.6cm}p{3.4cm}p{4.4cm}@{}}
      \toprule
      \textbf{Se busca} & \textbf{Operación} & \textbf{Notación} \\
      \midrule
      El resultado: $2^{3} = ?$ & Potenciación & $8$ \\
      La base: $?^{3} = 8$ & \azul{Radicación} & $\sqrt[3]{8} = 2$ \\
      El exponente: $2^{?} = 8$ & \destacar{Logaritmación} &
        $\log_{2}8 = 3$ \\
      \bottomrule
    \end{tabular}
    \end{center}
  \end{block}

  \vspace{0.15cm}
  \begin{alertblock}{Por eso hay dos operaciones inversas y no una}
    \scriptsize
    Porque en $a^{b}$ la base y el exponente \destacar{no son
    intercambiables}: $2^{3} = 8$ pero $3^{2} = 9$. La suma y el producto sí
    son conmutativos, y por eso les basta con una sola inversa.
  \end{alertblock}
\end{frame}

\seccion{2.4 Logaritmos}

\begin{frame}{Definición}
  \scriptsize
  \[ \boxed{\ \log_{b} a = c \quad\Longleftrightarrow\quad b^{c} = a\ } \]

  \vspace{0.08cm}
  \begin{columns}[T,onlytextwidth]
    \begin{column}{0.48\textwidth}
      \begin{block}{Cómo se lee}
        \tiny
        ``El logaritmo de $a$ en base $b$ es $c$'' equivale a
        \destacar{``a qué potencia hay que elevar $b$ para obtener $a$''}.

        \vspace{0.08cm}
        $\log_{2}8 = 3$ porque $2^{3}=8$.
      \end{block}
    \end{column}
    \begin{column}{0.48\textwidth}
      \begin{alertblock}{Las tres condiciones}
        \tiny
        \textbf{Base:} $b > 0$ y $b \neq 1$. \\
        \textbf{Argumento:} $a > 0$.

        \vspace{0.06cm}
        \destacar{No existe el logaritmo de cero ni de un número negativo}: no
        hay exponente que haga $2^{c} = -8$, porque una potencia de base
        positiva siempre es positiva.
      \end{alertblock}
    \end{column}
  \end{columns}

  \vspace{0.1cm}
  \begin{exampleblock}{Tres valores que salen de la definición, sin calcular}
    \scriptsize
    \[ \log_{b} 1 = 0 \quad (b^{0}=1) \qquad
       \log_{b} b = 1 \quad (b^{1}=b) \qquad
       \log_{b} b^{n} = n \]
  \end{exampleblock}
\end{frame}

\begin{frame}{Pasar de una forma a la otra}
  \scriptsize
  \begin{center}
  \begin{tabular}{@{}p{4.2cm}p{4.2cm}p{3.6cm}@{}}
    \toprule
    \textbf{Forma exponencial} & \textbf{Forma logarítmica} &
    \textbf{Se lee} \\
    \midrule
    $3^{4} = 81$ & $\log_{3}81 = 4$ & Base 3, exponente 4 \\
    $5^{-2} = \tfrac{1}{25}$ & $\log_{5}\tfrac{1}{25} = -2$ &
      Argumento menor que 1 $\Rightarrow$ log negativo \\
    $10^{3} = 1\,000$ & $\log 1\,000 = 3$ & Base 10: se omite \\
    $e^{4}$ & $\ln e^{4} = 4$ & Base $e$: se escribe $\ln$ \\
    $2^{10} = 1\,024$ & $\log_{2}1\,024 = 10$ & Base 2 \\
    \bottomrule
  \end{tabular}
  \end{center}

  \vspace{0.1cm}
  \begin{alertblock}{El truco para no equivocarse nunca}
    \scriptsize
    \destacar{La base del logaritmo es la base de la potencia, y el resultado
    del logaritmo es el exponente.} Dibuje una flecha desde la base hasta el
    resultado: $\log_{\azul{3}}81 = \destacar{4}$ se convierte en
    $\azul{3}^{\destacar{4}} = 81$.
  \end{alertblock}
\end{frame}

\begin{frame}{Las tres bases que se usan}
  \scriptsize
  \begin{center}
  \begin{tabular}{@{}p{2.4cm}p{2.0cm}p{3.0cm}p{4.4cm}@{}}
    \toprule
    \textbf{Nombre} & \textbf{Notación} & \textbf{Base} &
    \textbf{Dónde se usa} \\
    \midrule
    Decimal o común & $\log a$ & $10$ &
      pH, decibeles, escala Richter, órdenes de magnitud \\
    \addlinespace
    Natural & $\ln a$ & $e \approx 2{,}71828$ &
      Crecimiento continuo, cálculo, interés compuesto \\
    \addlinespace
    Binario & $\log_{2} a$ & $2$ &
      \destacar{Complejidad de algoritmos}, bits, árboles \\
    \bottomrule
  \end{tabular}
  \end{center}

  \vspace{0.1cm}
  \begin{exampleblock}{Por qué a un ingeniero de software le importa $\log_2$}
    \scriptsize
    Buscar un dato entre \azul{un millón} con búsqueda binaria toma
    $\log_{2}(1\,000\,000) \approx 19{,}9$, es decir \destacar{20 pasos}.
    Recorriendo uno por uno serían un millón. \azul{Esa diferencia es el
    logaritmo.}
  \end{exampleblock}

  \vspace{0.05cm}
  {\tiny\color{uniGris} La calculadora solo trae \fbox{log} y \fbox{ln}. Para
  cualquier otra base se usa el \destacar{cambio de base} que veremos en la
  sesión 11.}
\end{frame}

\begin{frame}{Ejemplo guiado — calcular logaritmos sin calculadora}
  \scriptsize
  \begin{center}
  \begin{tabular}{@{}lp{4.6cm}l@{}}
    \toprule
    \textbf{Logaritmo} & \textbf{Se pregunta} & \textbf{Resultado} \\
    \midrule
    $\log_{2}64$ & ¿$2$ a la cuánto da $64$? \ $2^{6}=64$ & $6$ \\
    $\log_{3}81$ & $3^{4}=81$ & $4$ \\
    $\log_{5}\tfrac{1}{25}$ & $5^{-2}=\tfrac{1}{25}$ & $-2$ \\
    $\log 0{,}001$ & $10^{-3}=0{,}001$ & $-3$ \\
    $\log_{4}2$ & $4^{1/2}=2$ & $\tfrac{1}{2}$ \\
    $\log_{1/2}8$ & $\left(\tfrac{1}{2}\right)^{-3}=8$ & $-3$ \\
    $\log_{7}1$ & $7^{0}=1$ & $0$ \\
    \bottomrule
  \end{tabular}
  \end{center}

  \vspace{0.08cm}
  \begin{alertblock}{El método, siempre el mismo}
    \tiny
    \destacar{Escriba el argumento como potencia de la base.} Si logra hacerlo,
    el logaritmo es ese exponente y no hace falta calculadora. Los tres casos
    difíciles son los de arriba: \azul{argumento fraccionario} (log negativo),
    \azul{argumento menor que la base} (log fraccionario) y \azul{base
    fraccionaria} (invierte el signo).
  \end{alertblock}
\end{frame}

\begin{frame}{Ejercicio 1 — en clase}
  \begin{alertblock}{Individual (12 minutos)}
    \scriptsize
    \textbf{Parte A.} Calcule sin calculadora:

    \vspace{0.15cm}
    \centering
    \begin{tabular}{@{}llll@{}}
      1. \ $\log_{2}32$ & 2. \ $\log_{3}\dfrac{1}{9}$ &
      3. \ $\log 10\,000$ & 4. \ $\log_{9}3$ \\[0.9em]
      5. \ $\ln e^{7}$ & 6. \ $\log_{6}1$ &
      7. \ $\log_{2}\dfrac{1}{16}$ & 8. \ $\log_{5}5$ \\
    \end{tabular}

    \vspace{0.2cm}
    \raggedright
    \textbf{Parte B.} Escriba en la otra forma:

    \vspace{0.1cm}
    \begin{columns}[T,onlytextwidth]
      \begin{column}{0.48\textwidth}
        \begin{enumerate}\setlength{\itemsep}{-0.02cm}
          \item $4^{3} = 64$
          \item $10^{-2} = 0{,}01$
        \end{enumerate}
      \end{column}
      \begin{column}{0.48\textwidth}
        \begin{enumerate}\setlength{\itemsep}{-0.02cm}
          \setcounter{enumi}{2}
          \item $\log_{7}49 = 2$
          \item $\ln x = 5$
        \end{enumerate}
      \end{column}
    \end{columns}
  \end{alertblock}
\end{frame}

\begin{frame}{Ejercicio 1 — solución}
  \scriptsize
  \begin{columns}[T,onlytextwidth]
    \begin{column}{0.49\textwidth}
      \begin{block}{Parte A}
        \tiny
        \begin{center}
        \begin{tabular}{@{}llll@{}}
          1 & $5$ & 2 & $-2$ \\
          3 & $4$ & 4 & $\tfrac{1}{2}$ \\
          5 & $7$ & 6 & $0$ \\
          7 & $-4$ & 8 & $1$ \\
        \end{tabular}
        \end{center}
      \end{block}
    \end{column}
    \begin{column}{0.49\textwidth}
      \begin{block}{Parte B}
        \tiny
        \textbf{1.} $\log_{4}64 = 3$

        \textbf{2.} $\log 0{,}01 = -2$

        \textbf{3.} $7^{2} = 49$

        \textbf{4.} $e^{5} = x$
      \end{block}
    \end{column}
  \end{columns}

  \vspace{0.1cm}
  \begin{alertblock}{Los dos que hay que revisar}
    \tiny
    \textbf{A4.} $\log_{9}3 = \tfrac{1}{2}$ porque $9^{1/2}=3$: cuando el
    argumento es \azul{menor} que la base, el logaritmo está entre 0 y 1. \quad
    \textbf{A7.} $\log_{2}\tfrac{1}{16} = -4$: argumento \azul{menor que 1}
    siempre da logaritmo \destacar{negativo}.
  \end{alertblock}
\end{frame}

\begin{frame}{Ecuaciones exponenciales sencillas}
  \scriptsize
  \begin{columns}[T,onlytextwidth]
    \begin{column}{0.48\textwidth}
      \begin{block}{Método 1 — igualar bases}
        \tiny
        Si se logra escribir los dos lados con la \azul{misma base}, los
        exponentes son iguales.
        \[ 2^{x} = 32 \;\Rightarrow\; 2^{x} = 2^{5}
           \;\Rightarrow\; \destacar{x = 5} \]
        \[ 3^{x+1} = 27 \;\Rightarrow\; 3^{x+1}=3^{3}
           \;\Rightarrow\; x+1 = 3 \;\Rightarrow\; \destacar{x=2} \]
      \end{block}
    \end{column}
    \begin{column}{0.48\textwidth}
      \begin{block}{Método 2 — aplicar logaritmo}
        \tiny
        Cuando las bases no se pueden igualar:
        \[ 2^{x} = 50 \;\Rightarrow\; x = \log_{2}50 \]
        \[ x = \frac{\log 50}{\log 2} \approx \destacar{5{,}64} \]
        {\color{uniGris} El cambio de base se justifica en la sesión 11; hoy
        basta con saber que existe.}
      \end{block}
    \end{column}
  \end{columns}

  \vspace{0.1cm}
  \begin{alertblock}{Cuál usar}
    \scriptsize
    \destacar{Siempre se intenta primero igualar bases}: da un resultado exacto
    y sin calculadora. Solo si el argumento no es una potencia de la base
    ($50$ no lo es de $2$) se recurre al logaritmo.
  \end{alertblock}
\end{frame}

\begin{frame}{Ejercicio 2 — en clase}
  \begin{alertblock}{En parejas (15 minutos)}
    \scriptsize
    Resuelva igualando bases:

    \vspace{0.25cm}
    \begin{columns}[T,onlytextwidth]
      \begin{column}{0.48\textwidth}
        \begin{enumerate}\setlength{\itemsep}{0.16cm}
          \item $5^{x} = 125$
          \item $2^{x-3} = 16$
          \item $3^{2x} = 81$
        \end{enumerate}
      \end{column}
      \begin{column}{0.48\textwidth}
        \begin{enumerate}\setlength{\itemsep}{0.16cm}
          \setcounter{enumi}{3}
          \item $\left(\dfrac{1}{2}\right)^{x} = 8$
          \item $4^{x} = 2$
          \item $10^{x} = 0{,}001$
        \end{enumerate}
      \end{column}
    \end{columns}
  \end{alertblock}
\end{frame}

\begin{frame}{Ejercicio 2 — solución}
  \tiny
  \begin{columns}[T,onlytextwidth]
    \begin{column}{0.49\textwidth}
      \begin{block}{Puntos 1 a 3}
        \tiny
        \begin{align*}
          1)\;& 5^{x}=5^{3} \Rightarrow \destacar{x=3} \\[0.3em]
          2)\;& 2^{x-3}=2^{4} \Rightarrow x-3=4 \Rightarrow
                \destacar{x=7} \\[0.3em]
          3)\;& 3^{2x}=3^{4} \Rightarrow 2x=4 \Rightarrow \destacar{x=2}
        \end{align*}
      \end{block}
    \end{column}
    \begin{column}{0.49\textwidth}
      \begin{block}{Puntos 4 a 6}
        \tiny
        \begin{align*}
          4)\;& 2^{-x}=2^{3} \Rightarrow -x=3 \Rightarrow
                \destacar{x=-3} \\[0.3em]
          5)\;& 2^{2x}=2^{1} \Rightarrow 2x=1 \Rightarrow
                \destacar{x=\tfrac{1}{2}} \\[0.3em]
          6)\;& 10^{x}=10^{-3} \Rightarrow \destacar{x=-3}
        \end{align*}
      \end{block}
    \end{column}
  \end{columns}

  \vspace{0.1cm}
  \begin{alertblock}{Los dos trucos que aparecen}
    \tiny
    En el \destacar{4}, $\left(\tfrac{1}{2}\right)^{x} = 2^{-x}$: la base
    fraccionaria se vuelve exponente negativo. \quad
    En el \destacar{5}, $4 = 2^{2}$: hay que \azul{cambiar la base para que
    coincida con el otro lado}, y ahí aparece el exponente fraccionario.
  \end{alertblock}
\end{frame}

\begin{frame}{Ejercicio 3 — en clase (taller principal)}
  \begin{alertblock}{En grupos de 3 (20 minutos)}
    \tiny
    \setbeamerfont{itemize/enumerate body}{size=\tiny}
    \begin{enumerate}\setlength{\itemsep}{0.05cm}
      \item Calcule: \ $\log_{2}128$, \ $\log_{4}64$, \ $\log_{8}2$, \
            $\log_{3}\tfrac{1}{81}$.
      \item ¿Para qué valores de $x$ está definido $\log_{5}(x-3)$?
      \item Un arreglo tiene 4\,096 elementos. ¿Cuántos pasos toma una búsqueda
            binaria? (Calcule $\log_{2}4\,096$.)
      \item ¿Y si el arreglo tuviera el doble de elementos? ¿Cuántos pasos
            más?
      \item Una población de bacterias se duplica cada hora y empieza en 500.
            Escriba la expresión del número de bacterias tras $t$ horas, y
            plantee (sin resolver) la ecuación para saber cuándo llega a
            32\,000.
      \item Encuentre el error: \ $\log_{2}(-8) = -3$
    \end{enumerate}
  \end{alertblock}
\end{frame}

\begin{frame}{Ejercicio 3 — solución}
  \scriptsize
  \begin{columns}[T,onlytextwidth]
    \begin{column}{0.48\textwidth}
      \begin{block}{Puntos 1 a 4}
        \tiny
        \textbf{1.} $7$; \ $3$; \ $\tfrac{1}{3}$; \ $-4$

        \vspace{0.06cm}
        \textbf{2.} Se exige $x-3 > 0$, luego $\destacar{x > 3}$.

        \vspace{0.06cm}
        \textbf{3.} $\log_{2}4\,096 = \destacar{12}$ pasos, porque
        $2^{12}=4\,096$.

        \vspace{0.06cm}
        \textbf{4.} $8\,192 = 2^{13}$: \destacar{un solo paso más}. Duplicar
        los datos cuesta un paso: eso es el crecimiento logarítmico.
      \end{block}
    \end{column}
    \begin{column}{0.50\textwidth}
      \begin{alertblock}{Puntos 5 y 6}
        \tiny
        \textbf{5.} $N(t) = 500\cdot 2^{t}$. Para llegar a 32\,000:
        \[ 500\cdot 2^{t} = 32\,000 \;\Rightarrow\; 2^{t} = 64 \]
        (y de ahí $t = 6$ horas).

        \vspace{0.08cm}
        \textbf{6.} \destacar{No existe}: el argumento de un logaritmo debe ser
        positivo. Lo que sí es cierto es $\log_{2}\tfrac{1}{8} = -3$.
        \azul{Un logaritmo puede ser negativo; su argumento, nunca.}
      \end{alertblock}
    \end{column}
  \end{columns}

  \vspace{0.05cm}
  {\tiny\color{uniGris} El punto \destacar{4} es la razón por la que la
  búsqueda binaria escala: pasar de mil a un millón de datos sube de 10 a 20
  pasos, no de mil a un millón.}
\end{frame}

\begin{frame}{Errores frecuentes en esta sesión}
  \scriptsize
  \begin{center}
  \begin{tabular}{@{}p{5.2cm}p{6.6cm}@{}}
    \toprule
    \textbf{Error} & \textbf{Corrección} \\
    \midrule
    Calcular $\log_{2}(-8)$ &
      No existe: el argumento debe ser \azul{positivo} \\
    \addlinespace
    Creer que un logaritmo no puede ser negativo &
      Sí puede: $\log_{2}\tfrac{1}{8} = -3$ \\
    \addlinespace
    Confundir base con argumento &
      En $\log_{3}81$, la base es 3 y el argumento 81 \\
    \addlinespace
    Escribir $\log 100 = 10$ &
      Es $2$: $10^{2}=100$. El log \azul{es el exponente} \\
    \addlinespace
    Usar base 1 &
      $1^{c}=1$ siempre: no se puede despejar nada \\
    \addlinespace
    Olvidar el dominio al resolver $\log(x-3)$ &
      Hay que exigir $x-3>0$ antes de operar \\
    \bottomrule
  \end{tabular}
  \end{center}
\end{frame}

\begin{frame}{Cierre de la sesión}
  \begin{block}{Lo que debe quedar claro hoy}
    \scriptsize
    \begin{itemize}\setlength{\itemsep}{0.02cm}
      \item $\log_{b}a = c \Leftrightarrow b^{c}=a$: \destacar{el logaritmo es
            el exponente}.
      \item Base $b>0$, $b\neq 1$; argumento $a>0$.
      \item $\log_{b}1 = 0$, \ $\log_{b}b = 1$, \ $\log_{b}b^{n} = n$.
      \item Tres bases usuales: $\log$ (10), $\ln$ ($e$) y $\log_2$
            (informática).
      \item Ecuaciones exponenciales: primero \azul{igualar bases}; si no se
            puede, aplicar logaritmo.
      \item Un logaritmo \azul{puede} ser negativo; su argumento \destacar{no}.
    \end{itemize}
  \end{block}

  \vspace{0.08cm}
  \begin{alertblock}{Trabajo independiente}
    \tiny
    Resolver los 20 ejercicios de la lámina siguiente. En la sesión 11 vemos
    las \destacar{propiedades de los logaritmos} —producto, cociente, potencia
    y cambio de base—, y hacemos el repaso y el simulacro del Examen 2.
  \end{alertblock}
\end{frame}

\begin{frame}{Ejercicios propuestos}
  \vspace{-0.1cm}
  \begin{beamercolorbox}[wd=\textwidth,sep=0.10cm]{block body}
    \centering\scriptsize
    \textbf{1--8:} calcular. \quad \textbf{9--14:} cambiar de forma. \quad
    \textbf{15--20:} ecuaciones y aplicaciones.
  \end{beamercolorbox}
  \listatiny
  \vspace{0.15cm}
  \begin{columns}[T,onlytextwidth]
    \begin{column}{0.47\textwidth}
      \begin{enumerate}\setlength{\itemsep}{0.16cm}
        \item $\log_{2}256$
        \item $\log_{3}243$
        \item $\log_{10}100\,000$
        \item $\log_{5}\dfrac{1}{125}$
        \item $\log_{16}4$
        \item $\log_{8}1$
        \item $\ln e^{9}$
        \item $\log_{1/3}9$
        \item Escriba $\log_{2}64 = 6$ en forma exponencial
        \item Escriba $6^{3}=216$ en forma logarítmica
      \end{enumerate}
    \end{column}
    \begin{column}{0.50\textwidth}
      \begin{enumerate}\setlength{\itemsep}{0.16cm}
        \setcounter{enumi}{10}
        \item Escriba $10^{-4}=0{,}0001$ en forma logarítmica
        \item Escriba $\ln x = 3$ en forma exponencial
        \item ¿Existe $\log_{5}0$?
        \item ¿Para qué $x$ existe $\log_{2}(x+4)$?
        \item Resuelva $3^{x}=243$
        \item Resuelva $2^{x+1}=64$
        \item Resuelva $\left(\tfrac{1}{5}\right)^{x}=25$
        \item Resuelva $9^{x}=3$
        \item ¿Cuántos pasos toma una búsqueda binaria en 1\,024 datos?
        \item ¿Por qué la base no puede ser 1?
      \end{enumerate}
    \end{column}
  \end{columns}
\end{frame}

\begin{frame}{Respuestas}{Autoevaluación de los ejercicios propuestos}
  \listatiny
  \vspace{0.15cm}
  \begin{columns}[T,onlytextwidth]
    \begin{column}{0.47\textwidth}
      \begin{enumerate}\setlength{\itemsep}{0.16cm}
        \item $8$
        \item $5$
        \item $5$
        \item $-3$
        \item $\tfrac{1}{2}$
        \item $0$
        \item $9$
        \item $-2$
        \item $2^{6}=64$
        \item $\log_{6}216 = 3$
      \end{enumerate}
    \end{column}
    \begin{column}{0.50\textwidth}
      \begin{enumerate}\setlength{\itemsep}{0.16cm}
        \setcounter{enumi}{10}
        \item $\log 0{,}0001 = -4$
        \item $e^{3}=x$
        \item No: el argumento debe ser positivo
        \item $x > -4$
        \item $x = 5$
        \item $x = 5$
        \item $x = -2$
        \item $x = \tfrac{1}{2}$
        \item $\log_{2}1\,024 = 10$ pasos
        \item Porque $1^{c}=1$ para todo $c$: no habría un único exponente
      \end{enumerate}
    \end{column}
  \end{columns}
\end{frame}

\begin{frame}{Referencias}
  \footnotesize
  \begin{itemize}
    \item[{[1]}] Ramos del Olmo, Á. M. (2016). \textit{Matemáticas básicas
      para el acceso a la universidad}. Ediciones Pirámide.
      \texttt{https://elibro.net/es/ereader/uniminuto/49134}
    \item[{[2]}] González Pulido, J. W., Chacón Benavides, J. A. \&
      Fonseca Correa, L. Á. (2023). \textit{Matemática básica} (1.\textsuperscript{a} ed.).
      Universidad Mariana.
      \texttt{https://elibro.net/es/ereader/uniminuto/280647}
    \item[{[3]}] Placencia Valero, J. (2008). \textit{Compendio de matemática
      básica elemental} (1.\textsuperscript{a} ed.). Editorial Tébar Flores.
      \texttt{https://elibro.net/es/ereader/uniminuto/269758}
  \end{itemize}

  \vspace{0.35cm}
  \begin{center}
    {\Large\color{uniAzul}\textbf{¿Preguntas?}}
  \end{center}
\end{frame}

\end{document}
